\section{民商事的法律适用}
\begin{enumerate}
  \item 占据的分值比例是非常的大的，多的时候多达 60\% 到 70\% 的知识点。
  \item 原则性规定。\qd{两个原则一句话}。\zd{意思自治、最密切联系原则}。
        \begin{enumerate}
          \item 意思自治。考试重复率比较高。有授权选择才有效。
                \begin{enumerate}
                  \item 授权突破实际联系。只有\zd{法律允许选择} + \zd{没有限制选择的范围}才有效（\zd{授权性规定}）。
                  \item \zd{一审法庭辩论终结前 + 援引相同 + 没有异议}。三缺一不可。
                  \item 中国\zd{已加入未加入}国际条约\zd{都可以}，违反公序良俗除外。
                  \item 口头书面均可。
                \end{enumerate}
          \item 最密切联系原则。兜底性原则。兜法律滞后、不完备的底。
          \item 涉外法律适用的规则\zd{参照适用}于解决涉港澳台纠纷（\zd{不叫涉外纠纷}）的法律适用。
        \end{enumerate}
  \item 民事主体。
        \begin{kk}{自然人}
          \begin{enumerate}
            \item \zd{自然人}内生性问题，\zd{自己的问题}，有别于与别人发生关系产生的纠纷。
            \item \zd{自然人用经居地}。民事权利能力、民事行为能力、宣告失踪、宣告死亡、人格权内容等。
            \item \zd{居无行有行为地}。民事行为能力适用其本国，从事民事活动适用行为地。
          \end{enumerate}
        \end{kk}

        \begin{kk}{法人}
          \zd{法人登记营业地，排除机构所在地}。法人内生性问题。股东权利、义务、组织机构等。
        \end{kk}

        \begin{kk}{时效、代理、信托}
          \begin{enumerate}
            \item 时效。买卖合同时效，买卖合同为基础法。\zd{时效基础法一致}。从属性。
            \item 代理。委托代理允许意思自治优先，指定/法定代理无效。
            \item 信托。财产所在地 \hz 关系发生地。\al{甲国乙与丙信托公司在北京签订协议}。
          \end{enumerate}
          \tcblower
          \zd{委代信托可自治，代理内部关系地，代理外部行为地，信托财产或关系}。
        \end{kk}
\end{enumerate}

